\documentclass[12pt]{article}
\usepackage[left = 0.6in, right = 0.6in, top = 0.4in, bottom = 1.5in]{geometry} 
\usepackage{pgfplots, fontawesome5, booktabs, xcolor, hyperref, amsmath,amsthm,amssymb,amsfonts, enumitem, fancyhdr, color, comment, graphicx, environ, actuarialsymbol}
\pagestyle{fancy}
\setlength{\headheight}{65pt}
%%%%%%%%%%%%%%%%%%%%%%%%%%%%%%%%%%%%%%%%%%%%%
\lhead{ }
\rhead{Math Methods 1/2 review}
%%%%%%%%%%%%%%%%%%%%%%%%%%%%%%%%%%%%%%%%%%%%%
\newcommand{\emptygrid}[3]{ % #1: range, #2: y-label, #3: "labels" or "nolabels"
    \begin{tikzpicture}
    \begin{axis}[
        axis lines = middle,
        xlabel = $x$,
        ylabel = $#2$,
        xmin=-#1, xmax=#1,
        ymin=-#1, ymax=#1,
        grid = both,
        grid style = {line width=.1pt, draw=gray!30},
        xtick distance = 5,
        ytick distance = 5,
        minor tick num = 4,
        tick label style = {font=\tiny},
        width=10cm, height=10cm,
        axis equal image,
        % Logic to hide labels if #3 is "nolabels"
        xticklabels={ \ifx\empty#3\empty\else\ifnum\pdfstrcmp{#3}{nolabels}=0 \else \fi\fi },
        yticklabels={ \ifx\empty#3\empty\else\ifnum\pdfstrcmp{#3}{nolabels}=0 \else \fi\fi }
    ]
    \end{axis}
    \end{tikzpicture}
}
%%%%%%%%%%%%%%%%%%%%%%%%%%%%%%%%%%%%%%%%%%%%%
\newtheorem{theorem}{Theorem}[section]
\begin{document}
Don't worry, this is not a real test! This is just a collection of problems that lets me know what you know. If you dont know how to do a question, feel free to leave it blank or write down initial thoughts
\begin{itemize}
    \item You can use any reference you want
    \item only use a calculator on questions with (\faCalculator)  on them
    \item keep track of the time you spend on each question
\end{itemize}

\hrulefill

\begin{enumerate}
    \item Set and interval notation 
    \begin{enumerate}
        \item Translate into interval notation: "real numbers greater than -2 \& less than or equal to 5"
        \vspace{1cm}
        \item Translate into set notation: "integers greater than -2 \& less than or equal to 5"
        \vspace{1cm}
        \item Translate into interval notation: "real numbers {less than or equal to} -2 or {bigger than} 5"
        \vspace{1cm}
        \item Translate into set notation: $(-\infty, 3) \cup (3, \infty)$
        \vspace{1cm}
        \item Try translating (c) into set notation (A bit trickier)
        \vspace{1cm}
    \end{enumerate}

    \item Basic pseudocode: What value does $x$ have after executing the following pseudocode?
    \begin{verbatim}
    x <- 0 # this means let x take the value 0
    while x < 10 do
        if x is even then
            x <- x + 3 # remember that 0 is even :)
        else
            x <- x + 1
    \end{verbatim}
    \newpage

    \item Let $f: \mathbb{R} \to \mathbb{R}, f(x) = x^{3}$ and $g: \mathbb{R} \to \mathbb{R}, g(x) = 2(x-3)^{3} - 3$
    \begin{enumerate}
        \item Find a series of transformations that takes $f$ to $g$
        \vspace{5cm}
        \item Are there any points of intersection between $f$ and $g$? If so, find them. (\faCalculator)
        \vspace{5cm}
        \item Plot the graph of $h(x) = \frac{f(x)}{x} + 2x$ and label key points
        \begin{center}
            \emptygrid{10}{h(x)}{labels}
        \end{center}
        \newpage
        \item Find the slope of a tangent line to $g$ at $x = 3$ and $x = 4$
        \vspace{5cm}
        \item Find the \textit{signed} area under the curve $g$ from $x = 3$ to $x = 6$ 
        \vspace{5cm}
        \item Find the \textit{unsigned} area between the curve $g$ and the $x$ axis from $x = 3$ to $x = 6$ 
        \vspace{5cm}
    \end{enumerate}



    \item Let $f: [0, \frac{13\pi}{12}) \to \mathbb{R}, f(x) = \sin(2x) + 3$
    \begin{enumerate}
        \item Find the period and amplitude of $f$
        \vspace{1cm}
        \item Find the endpoints of the graph of $f$
        \vspace{3cm}
        \item Sketch the graph of $f$ over its domain, labelling key points
        \begin{center}
            \emptygrid{8}{f(x)}{nolabels}
        \end{center}
    \end{enumerate}

    \item You have a sheet of paper with thickness of $0.1mm$. You keep folding it in half, and measure the thickness after each fold.
    \begin{enumerate}
        \item What is the thickness after 3 folds?
        \vspace{2cm}
        \item let $f(n)$ be the thickness in mm after $n$ folds. Write the function $f$ in terms of $n$
        \vspace{3cm}
        \item The moon is $384400 km = 384400 \times 10^{6} mm$ away from the surface of earth. How many folds do you need for the paper to reach the moon? (\faCalculator)
        \vspace{4cm}
    \end{enumerate}

    \item Every morning before school, you pick a pair of socks at random from your drawer without looking. Suppose you have 3 pink, 8 white, and 7 black socks in the drawer.
    \begin{enumerate}
        \item What is the probability of picking a pair of matching white socks? (Hint: try a tree diagram)
        \vspace{3cm}
        \item What is the probability of picking a pair of matching socks (any colour)?
        \vspace{3cm}
        \item Given you chose a pair of matching socks, what is the probability that you chose a pair of black socks?
        \vspace{3cm}
    \end{enumerate}
    \item Counting: Now dont worry if you dont know how to do these, there is a lot of variability in how different schools cover this topic (if they cover it at all) but having a decent understanding of these concepts can be extremely useful in probability/stats topics later on.
    \begin{enumerate}
    \item How many different ways can you arrange the letters in the word "METHODS"? (\faCalculator)
    \vspace{3cm}
    \item Explain what ${10 \choose{4}}$ means and how do you calculate it? Explain why the formula for ${n \choose{k}}$ makes sense
    \vspace{3cm}
    
\end{enumerate}
    
\end{enumerate}



\end{document}